%!TEX root = ../main.tex
% =============================================================
%  CAPÍTULO 1 — INTRODUCCIÓN            (Entrega 1 · 5 % · 9/6)
% =============================================================

\chapter{Introducción}\label{cap:introduccion}

\begin{guia}[Guía de la cátedra: qué se evalúa en este capítulo]
\begin{enumerate}
    \item \textbf{Descripción del problema}: concreto, anclado en evidencia
          (datos duros con fuente, casos reales, competidores nombrados).
    \item \textbf{Objetivo general}: uno solo, medible, en un párrafo.
          Los \textbf{objetivos específicos} se desprenden de él y
          \emph{no} son los pasos del plan de trabajo.
    \item \textbf{Hipótesis}: explícita, formulada como pregunta o
          afirmación de investigación que pueda responderse en el
          capítulo~\ref{cap:conclusiones}.
    \item \textbf{Relevancia}: por qué importa, a quién beneficia, qué les
          falta a las soluciones existentes.
    \item \textbf{Alcance y recursos}: qué entra y qué no; stack y fuentes
          de datos previstos. El ``alcance que nunca cierra'' es el riesgo
          número uno de la materia.
\end{enumerate}
Extensión orientativa: 4 a 6 páginas. Redacción impersonal, sin ``yo
creo'', oraciones de 15 a 22 palabras, sin palabras vacías
(``claramente'', ``obviamente''). Toda cifra con su fuente citada, por
ejemplo \citep{breiman2001}.
\end{guia}

\completar{Desarrollo del capítulo. Las secciones y subsecciones se
organizan como resulte mejor para el proyecto.}

\section{\completar{Nombre de la sección}}

\begin{ejemplo}[Ejemplo: forma de un objetivo general]
Desarrollar un sistema de \completar{qué} para \completar{quién} en
\completar{dónde}, que permita \completar{resultado medible} mediante
\completar{enfoque o tecnología central}.
\end{ejemplo}

\completar{Texto.}

\subsection{\completar{Nombre de la subsección}}

\completar{Texto.}
